\textbf{Phần 1. Trắc nghiệm lựa chọn}
\setcounterpageref{ex}{0}

\begin{ex} %[5H1V8]
Một người thợ xây dùng các viên gạch men hình vuông cạnh $40\text{ cm}$ để lát một căn phòng hình chữ nhật có chiều dài $8\text{ m}$ và chiều rộng $6\text{ m}$. Coi diện tích mạch vữa là không đáng kể. Số viên gạch men cần dùng để lát kín căn phòng là
\choice
{$200$ viên}
{\True $300$ viên}
{$400$ viên}
{$250$ viên}
\loigiai{
Diện tích căn phòng là:
\[
8 \times 6 = 48\text{ (m}^2\text{)} = 480\,000\text{ cm}^2.
\]
Diện tích một viên gạch men hình vuông là:
\[
40 \times 40 = 1\,600\text{ (cm}^2\text{)}.
\]
Số viên gạch men cần dùng là:
\[
480\,000 : 1\,600 = 300\text{ (viên)}.
\]
\textbf{Đáp án của câu hỏi là B}
}
\end{ex}

\begin{ex} %[5H1V8]
Một khu đất hình chữ nhật có chu vi $120\text{ m}$. Nếu tăng chiều rộng thêm $5\text{ m}$ và giảm chiều dài đi $5\text{ m}$ thì khu đất trở thành hình vuông. Diện tích ban đầu của khu đất đó là
\choice
{$800\text{ m}^2$}
{\True $875\text{ m}^2$}
{$900\text{ m}^2$}
{$850\text{ m}^2$}
\loigiai{
Nửa chu vi khu đất là:
\[
120 : 2 = 60\text{ (m)}.
\]
Khi tăng chiều rộng $5\text{ m}$ và giảm chiều dài $5\text{ m}$ thì hai cạnh bằng nhau, do đó chiều dài hơn chiều rộng là:
\[
5 + 5 = 10\text{ (m)}.
\]
Chiều dài khu đất ban đầu là:
\[
(60 + 10) : 2 = 35\text{ (m)}.
\]
Chiều rộng khu đất ban đầu là:
\[
60 - 35 = 25\text{ (m)}.
\]
Diện tích ban đầu của khu đất là:
\[
35 \times 25 = 875\text{ (m}^2\text{)}.
\]
\textbf{Đáp án của câu hỏi là B}
}
\end{ex}

\begin{ex} %[5H1V8]
Lúc $4$ giờ đúng, hai kim giờ và kim phút của một chiếc đồng hồ tạo thành một góc có số đo là bao nhiêu và là góc gì?
\choice
{$90^\circ$ -- Góc vuông}
{\True $120^\circ$ -- Góc tù}
{$150^\circ$ -- Góc tù}
{$60^\circ$ -- Góc nhọn}
\loigiai{
Mặt đồng hồ được chia làm $12$ khoảng bằng nhau, mỗi khoảng tương ứng với số đo góc là:
\[
360^\circ : 12 = 30^\circ.
\]
Lúc $4$ giờ đúng, kim dài chỉ số $12$, kim ngắn chỉ số $4$. Khoảng cách giữa hai kim là $4$ khoảng.\\
Số đo góc giữa hai kim là:
\[
4 \times 30^\circ = 120^\circ.
\]
Vì $90^\circ < 120^\circ < 180^\circ$ nên đây là góc tù.\\
\textbf{Đáp án của câu hỏi là B}
}
\end{ex}

\begin{ex} %[5H1V8]
Một thửa ruộng hình bình hành có độ dài đáy là $45\text{ m}$, chiều cao bằng $\dfrac{2}{3}$ độ dài đáy. Cứ mỗi $100\text{ m}^2$ người ta thu hoạch được $60\text{ kg}$ thóc. Hỏi cả thửa ruộng thu hoạch được bao nhiêu tạ thóc?
\choice
{$81$ tạ}
{\True $8{,}1$ tạ}
{$810$ tạ}
{$8\,100$ tạ}
\loigiai{
Chiều cao của thửa ruộng là:
\[
45 \times \dfrac{2}{3} = 30\text{ (m)}.
\]
Diện tích thửa ruộng hình bình hành là:
\[
45 \times 30 = 1\,350\text{ (m}^2\text{)}.
\]
Số thóc thu hoạch được là:
\[
(1\,350 : 100) \times 60 = 810\text{ (kg)}.
\]
Đổi $810\text{ kg} = 8{,}1$ tạ.\\
\textbf{Đáp án của câu hỏi là B}
}
\end{ex}

\begin{ex} %[5H1V8]
Một tấm bìa hình thoi có độ dài hai đường chéo lần lượt là $2\text{ dm } 4\text{ cm}$ và $1\text{ dm } 5\text{ cm}$. Diện tích của tấm bìa đó là
\choice
{$360\text{ cm}^2$}
{\True $180\text{ cm}^2$}
{$18\text{ cm}^2$}
{$720\text{ cm}^2$}
\loigiai{
Đổi đơn vị: $2\text{ dm } 4\text{ cm} = 24\text{ cm}$; $1\text{ dm } 5\text{ cm} = 15\text{ cm}$.\\
Diện tích tấm bìa hình thoi là:
\[
\dfrac{24 \times 15}{2} = 180\text{ (cm}^2\text{)}.
\]
\textbf{Đáp án của câu hỏi là B}
}
\end{ex}

\begin{ex} %[5H1V8]
Người ta rào xung quanh một khu vườn hình chữ nhật có chiều dài $28\text{ m}$, chiều rộng $16\text{ m}$. Người ta để một lối ra vào rộng $2\text{ m}$. Hỏi chiều dài hàng rào bao quanh khu vườn là bao nhiêu mét?
\choice
{$88\text{ m}$}
{\True $86\text{ m}$}
{$84\text{ m}$}
{$90\text{ m}$}
\loigiai{
Chu vi của khu vườn hình chữ nhật là:
\[
(28 + 16) \times 2 = 88\text{ (m)}.
\]
Chiều dài hàng rào cần làm là:
\[
88 - 2 = 86\text{ (m)}.
\]
\textbf{Đáp án của câu hỏi là B}
}
\end{ex}

\begin{ex} %[5H1V8]
Một ô tô chở $3\text{ tấn } 5\text{ tạ}$ gạo và $4\text{ tạ } 50\text{ kg}$ ngô. Tổng khối lượng hàng ô tô đó đã chở theo đơn vị ki-lô-gam là
\choice
{$3\,550\text{ kg}$}
{\True $3\,950\text{ kg}$}
{$3\,905\text{ kg}$}
{$4\,000\text{ kg}$}
\loigiai{
Đổi sang đơn vị ki-lô-gam:
\begin{itemize}
    \item $3\text{ tấn } 5\text{ tạ} = 3\,000\text{ kg} + 500\text{ kg} = 3\,500\text{ kg}$.
    \item $4\text{ tạ } 50\text{ kg} = 400\text{ kg} + 50\text{ kg} = 450\text{ kg}$.
\end{itemize}
Tổng khối lượng hàng là:
\[
3\,500 + 450 = 3\,950\text{ (kg)}.
\]
\textbf{Đáp án của câu hỏi là B}
}
\end{ex}

\begin{ex} %[5H1V8]
Một vận động viên hoàn thành chặng đua thứ nhất trong $1\text{ giờ } 15\text{ phút}$ và chặng đua thứ hai trong $55\text{ phút}$. Tổng thời gian vận động viên hoàn thành cả hai chặng đua là
\choice
{$1\text{ giờ } 70\text{ phút}$}
{\True $2\text{ giờ } 10\text{ phút}$}
{$2\text{ giờ } 20\text{ phút}$}
{$2\text{ giờ } 5\text{ phút}$}
\loigiai{
Tổng thời gian là:
\[
1\text{ giờ } 15\text{ phút} + 55\text{ phút} = 1\text{ giờ } 70\text{ phút} = 2\text{ giờ } 10\text{ phút}.
\]
\textbf{Đáp án của câu hỏi là B}
}
\end{ex}

\begin{ex} %[5H1V8]
Hình thoi $ABCD$ có diện tích bằng $72\text{ cm}^2$, độ dài một đường chéo là $12\text{ cm}$. Độ dài đường chéo còn lại là
\choice
{$6\text{ cm}$}
{\True $12\text{ cm}$}
{$18\text{ cm}$}
{$24\text{ cm}$}
\loigiai{
Độ dài đường chéo còn lại là:
\[
72 \times 2 : 12 = 12\text{ (cm)}.
\]
\textbf{Đáp án của câu hỏi là B}
}
\end{ex}

\begin{ex} %[5H1V8]
Năm $1945$, Bác Hồ đọc Tuyên ngôn Độc lập khai sinh ra nước Việt Nam Dân chủ Cộng hòa. Năm đó thuộc thế kỉ nào?
\choice
{Thế kỉ XVIII}
{Thế kỉ XIX}
{\True Thế kỉ XX}
{Thế kỉ XXI}
\loigiai{
Năm $1945$ thuộc thế kỉ thứ $20$ (thế kỉ XX), kéo dài từ năm $1901$ đến năm $2000$.\\
\textbf{Đáp án của câu hỏi là C}
}
\end{ex}

\begin{ex} %[5H1V8]
Một cái sân hình vuông có cạnh $15\text{ m}$. Người ta mở rộng sân về cả bốn phía, mỗi phía thêm $2\text{ m}$. Diện tích phần sân tăng thêm là
\choice
{$136\text{ m}^2$}
{\True $136\text{ m}^2$}
{$120\text{ m}^2$}
{$144\text{ m}^2$}
\loigiai{
Cạnh của cái sân sau khi mở rộng là:
\[
15 + 2 + 2 = 19\text{ (m)}.
\]
Diện tích sân ban đầu là:
\[
15 \times 15 = 225\text{ (m}^2\text{)}.
\]
Diện tích sân mới sau khi mở rộng là:
\[
19 \times 19 = 361\text{ (m}^2\text{)}.
\]
Diện tích phần tăng thêm là:
\[
361 - 225 = 136\text{ (m}^2\text{)}.
\]
\textbf{Đáp án của câu hỏi là A}
}
\end{ex}

\begin{ex} %[5H1V8]
Có bốn bao xi măng, mỗi bao nặng $50\text{ kg}$. Hỏi cần thêm bao nhiêu bao xi măng như thế nữa để được tất cả $1\text{ tấn}$ xi măng?
\choice
{$12$ bao}
{\True $16$ bao}
{$18$ bao}
{$20$ bao}
\loigiai{
Đổi $1\text{ tấn} = 1\,000\text{ kg}$.\\
Bốn bao xi măng nặng:
\[
4 \times 50 = 200\text{ (kg)}.
\]
Khối lượng xi măng còn thiếu là:
\[
1\,000 - 200 = 800\text{ (kg)}.
\]
Số bao xi măng cần thêm là:
\[
800 : 50 = 16\text{ (bao)}.
\]
\textbf{Đáp án của câu hỏi là B}
}
\end{ex}

\textbf{Phần 2. Trắc nghiệm đúng sai}
\setcounterpageref{ex}{0}

\begin{ex} %[5H1V8]
Một mảnh đất hình chữ nhật có nửa chu vi là $45\text{ m}$, chiều rộng bằng $\dfrac{4}{5}$ chiều dài. Người ta dùng $\dfrac{1}{3}$ diện tích mảnh đất để đào ao thả cá. Xét tính đúng sai của các khẳng định sau:
\choiceTF
{\True Chiều dài của mảnh đất là $25\text{ m}$.}
{\True Chiều rộng của mảnh đất là $20\text{ m}$.}
{Diện tích cả mảnh đất là $450\text{ m}^2$.}
{\True Diện tích phần đất đào ao thả cá là $167\text{ m}^2$ (làm tròn đến hàng đơn vị).}
\loigiai{
\begin{itemize}
    \item Tổng số phần: $4 + 5 = 9$ phần.
    Chiều rộng: $45 : 9 \times 4 = 20\text{ (m)}$. Khẳng định b \textbf{Đúng}.
    Chiều dài: $45 - 20 = 25\text{ (m)}$. Khẳng định a \textbf{Đúng}.
    \item Diện tích mảnh đất: $25 \times 20 = 500\text{ (m}^2\text{)} 
e 450\text{ m}^2$. Khẳng định c \textbf{Sai}.
    \item Diện tích đào ao là: $500 \times \dfrac{1}{3} = \dfrac{500}{3} \approx 166{,}67 \approx 167\text{ (m}^2\text{)}$. Khẳng định d \textbf{Đúng}.
\end{itemize}
\begin{itemize}
    \item a) \textbf{Đúng}.
    \item b) \textbf{Đúng}.
    \item c) \textbf{Sai}: Diện tích là $500\text{ m}^2$.
    \item d) \textbf{Đúng}.
\end{itemize}
}
\end{ex}

\begin{ex} %[5H1V8]
Một xe tải chở hàng có lộ trình trong một ngày được ghi lại như sau:
Buổi sáng xe chạy $2\text{ giờ } 45\text{ phút}$. Buổi chiều xe chạy $3\text{ giờ } 30\text{ phút}$. Mỗi giờ xe chở hàng tiêu thụ $12\text{ lít}$ dầu. Xét tính đúng sai của các khẳng định sau:
\choiceTF
{\True Tổng thời gian xe chạy trong cả ngày là $6\text{ giờ } 15\text{ phút}$.}
{Thời gian xe chạy buổi chiều nhiều hơn buổi sáng là $35\text{ phút}$.}
{\True Đổi $6\text{ giờ } 15\text{ phút} = 6\dfrac{1}{4}\text{ giờ}$.}
{\True Cả ngày xe đã tiêu thụ hết $75\text{ lít}$ dầu.}
\loigiai{
\begin{itemize}
    \item Tổng thời gian: $2\text{ giờ } 45\text{ phút} + 3\text{ giờ } 30\text{ phút} = 5\text{ giờ } 75\text{ phút} = 6\text{ giờ } 15\text{ phút}$. Khẳng định a \textbf{Đúng}.
    \item Chiều hơn sáng: $3\text{ giờ } 30\text{ phút} - 2\text{ giờ } 45\text{ phút} = 45\text{ phút} 
e 35$ phút. Khẳng định b \textbf{Sai}.
    \item $15\text{ phút} = \dfrac{15}{60}\text{ giờ} = \dfrac{1}{4}\text{ giờ} \implies 6\text{ giờ } 15\text{ phút} = 6\dfrac{1}{4}\text{ giờ}$. Khẳng định c \textbf{Đúng}.
    \item Lượng dầu tiêu thụ: $12 \times 6\dfrac{1}{4} = 12 \times \dfrac{25}{4} = 75\text{ (lít)}$. Khẳng định d \textbf{Đúng}.
\end{itemize}
\begin{itemize}
    \item a) \textbf{Đúng}.
    \item b) \textbf{Sai}.
    \item c) \textbf{Đúng}.
    \item d) \textbf{Đúng}.
\end{itemize}
}
\end{ex}

\textbf{Phần 3. Trả lời ngắn}
\setcounterpageref{ex}{0}

\begin{ex} %[5H1V8]
Một thửa ruộng hình chữ nhật có chu vi $180\text{ m}$. Chiều dài hơn chiều rộng $20\text{ m}$. Hãy tính diện tích thửa ruộng đó theo đơn vị mét vuông.
\shortans{1925}
\loigiai{
\textbf{Đáp số:} $1925$\\
Nửa chu vi thửa ruộng là:
\[
180 : 2 = 90\text{ (m)}.
\]
Chiều dài thửa ruộng là:
\[
(90 + 20) : 2 = 55\text{ (m)}.
\]
Chiều rộng thửa ruộng là:
\[
90 - 55 = 35\text{ (m)}.
\]
Diện tích thửa ruộng là:
\[
55 \times 35 = 1\,925\text{ (m}^2\text{)}.
\]
}
\end{ex}

\begin{ex} %[5H1V8]
Một hình thoi có diện tích $216\text{ cm}^2$ và độ dài một đường chéo là $18\text{ cm}$. Hãy tính độ dài của đường chéo kia theo đơn vị xăng-ti-mét.
\shortans{24}
\loigiai{
\textbf{Đáp số:} $24$\\
Độ dài đường chéo còn lại là:
\[
216 \times 2 : 18 = 24\text{ (cm)}.
\]
}
\end{ex}

\begin{ex} %[5H1V8]
Có bao nhiêu phút trong $3\text{ giờ } 45\text{ phút}$?
\shortans{225}
\loigiai{
\textbf{Đáp số:} $225$\\
Ta có:
\[
3 \times 60 + 45 = 180 + 45 = 225\text{ (phút)}.
\]
}
\end{ex}

\begin{ex} %[5H1V8]
Một con voi nặng $4\text{ tấn } 800\text{ kg}$. Một con hà mã nặng bằng $\dfrac{5}{8}$ khối lượng con voi đó. Hỏi con hà mã nặng bao nhiêu tạ?
\shortans{30}
\loigiai{
\textbf{Đáp số:} $30$\\
Đổi $4\text{ tấn } 800\text{ kg} = 4\,800\text{ kg}$.\\
Khối lượng con hà mã là:
\[
4\,800 \times \dfrac{5}{8} = 3\,000\text{ (kg)}.
\]
Đổi $3\,000\text{ kg} = 30$ tạ.
}
\end{ex}

\begin{ex} %[5H1V8]
Một căn phòng hình chữ nhật có kích thước $9\text{ m} \times 4\text{ m}$. Người ta lát kín sàn bằng các viên gạch men hình vuông có diện tích $900\text{ cm}^2$. Hỏi cần bao nhiêu viên gạch men như thế?
\shortans{400}
\loigiai{
\textbf{Đáp số:} $400$\\
Diện tích căn phòng là:
\[
9 \times 4 = 36\text{ (m}^2\text{)} = 360\,000\text{ cm}^2.
\]
Số viên gạch men cần dùng là:
\[
360\,000 : 900 = 400\text{ (viên)}.
\]
}
\end{ex}

\begin{ex} %[5H1V8]
Một tấm tôn hình bình hành có diện tích $1\,200\text{ cm}^2$ và độ dài đáy là $40\text{ cm}$. Chiều cao tương ứng của tấm tôn đó bằng bao nhiêu xăng-ti-mét?
\shortans{30}
\loigiai{
\textbf{Đáp số:} $30$\\
Chiều cao của tấm tôn hình bình hành là:
\[
1\,200 : 40 = 30\text{ (cm)}.
\]
}
\end{ex}
